%% kapa: exam=1 round=36 year=2025 period=1 form=A subject=real_estate subject_ko=부동산학원론 seq=1 number=81
%% source: Q-Net 감정평가사 기출문제 (cst00309) artlSeq=5231525 file=raw/1008/5231525/2025년 제36회 감정평가사 1차시험 1교시 문제지 원본.pdf
%% conversion: kordoc 4.24.51 -> tools/gen1.py (auto); review=proofread-sonnet; flags=none
\begin{question}{81}
토지에 관한 설명으로 옳지 않은 것은?
\begin{choices}
  \item “토지의 표시”란 지적공부에 토지의 소재ㆍ지번(地番)ㆍ지목(地目)ㆍ면적ㆍ경계 또는 좌표를 등록한 것을 말한다.
  \item “지번”이란 필지에 부여하여 지적공부에 등록한 번호를 말한다.
  \item “토지의 이동(異動)”이란 홍수나 산사태 등으로 인해 토지의 지형이 변경된 것을 말한다.
  \item “합병”이란 지적공부에 등록된 2필지 이상을 1필지로 합하여 등록하는 것을 말한다.
  \item “분할”이란 지적공부에 등록된 1필지를 2필지 이상으로 나누어 등록하는 것을 말한다.
\end{choices}
\end{question}
\answer{3}
